% !TEX root = ../../../main.tex
\section{Linear combinations and bases}
\label{sec:linear-algebra-i:linear-combinations-and-bases}

% Sources: supplied lecture photos, quiz-prep pp. 3--5, and September 24
% handout, slide 20.
\subsection{Linear combinations and span}

\begin{definition}
  A linear combination of \(v_1,\ldots,v_n\in V\) is a vector of the form
  \[
    a_1v_1+\cdots+a_nv_n,\qquad a_1,\ldots,a_n\in\mathbb F.
  \]
  For \(S\subseteq V\), its span, denoted \(\operatorname{span}(S)\), is the
  set of all finite linear combinations of elements of \(S\).
\end{definition}

The empty linear combination is the zero vector, so
\(\operatorname{span}(\varnothing)=\{0\}\). A vector space is
\emph{finite-dimensional} when it has a finite spanning set. Unless explicitly
stated otherwise, the vector spaces in these notes are finite-dimensional.

\subsection{Linear independence}

\begin{definition}
  A subset \(S\subseteq V\) is linearly independent if, for every finite
  selection of distinct vectors \(v_1,\ldots,v_n\in S\),
  \[
    a_1v_1+\cdots+a_nv_n=0
    \quad\Longrightarrow\quad
    a_1=\cdots=a_n=0.
  \]
  The empty set is linearly independent.
\end{definition}

Equivalently, coefficients in a linear combination of distinct vectors from
an independent set are unique. If
\[
  a_1v_1+\cdots+a_nv_n=b_1v_1+\cdots+b_nv_n,
\]
then \(a_j=b_j\) for each \(j\).

\begin{definition}
  A subset is linearly dependent if some finite selection of distinct
  vectors from it satisfies
  \[
    a_1v_1+\cdots+a_nv_n=0
  \]
  with at least one coefficient nonzero.
\end{definition}
% Author review: the quiz-prep gloss about "maximal access" to the span and
% generating any vector in R^n from n vectors does not state a spanning
% hypothesis. The gloss about dependent vectors lying on the same line is
% also too restrictive. The precise supplied definitions are retained above.

\subsection{Bases and dimension}

\begin{definition}
  A basis of \(V\) is a linearly independent subset that spans \(V\).
  The dimension of a finite-dimensional vector space is the number of
  vectors in a basis.
\end{definition}

Every vector space has a basis. The zero vector space has the empty basis,
and a basis of \(\R^3\) has three vectors, so \(\dim\R^3=3\).

\subsection{Subspaces}

\begin{definition}
  A nonempty subset \(T\subseteq V\) is a subspace if it is closed under
  linear combinations. That is, for \(v,w\in T\) and \(a,b\in\mathbb F\),
  \[
    av+bw\in T.
  \]
  A subspace inherits addition and scalar multiplication from \(V\).
\end{definition}

Equivalently, every finite linear combination of elements of \(T\) lies in
\(T\). For finite-dimensional \(V\),
\[
  \dim T\leq\dim V,
  \qquad
  \dim T=\dim V\quad\Longleftrightarrow\quad T=V.
\]

\subsection{Bases beyond finite dimension}

In infinite-dimensional spaces, a basis \(S\) is still linearly independent, and
each vector is a \emph{finite} linear combination of elements of \(S\).

\begin{example}
  The vector space \(\mathcal P\) of polynomials in \(x\) with real
  coefficients has the basis
  \[
    \{1,x,x^2,\ldots\}.
  \]
\end{example}

For the vector space \(C(\R)\) of continuous functions \(\R\to\R\),
the existence of a basis follows from Zorn's lemma. No explicit basis is
constructed here. This example motivates the study of topological vector
spaces.
% The handout's stronger wording, "no explicit one can be written down",
% is not developed here into a separate impossibility claim or proof.
